\subsection*{Eigenvalues and Eigenvectors}
\paragraph*{Definition}
Let $A \in \mathbb{R}^{n\times n}$ and $\vec{x} \in \mathbb{R}^n$ be a nonzero vector. If $A\vec{x} = \lambda \vec{x}$ for some scalar $\lambda$, then
\begin{itemize}
    \item $\lambda$ is called an eigenvalue of $A$
    \item $\vec{x}$ is called an eigenvector of $A$ corresponding to $\lambda$
\end{itemize}

\paragraph*{Characteristic Polynomial}
The characteristic polynomial of $A$ is defined as $P_A(\lambda) = \det(A - \lambda I)$, where $I$ is the $n\times n$ identity matrix. 
The roots of $P_A(\lambda)$ are the eigenvalues of $A$. The eigenvectors corresponding to an eigenvalue $\lambda$ are the nonzero solutions 
to the homogeneous system $(A - \lambda I)\vec{x} = \vec{0}$. The set of all eigenvectors corresponding to $\lambda$ together with the zero 
vector is called the eigenspace of $A$ corresponding to $\lambda$.

\paragraph*{Theorem}
Let $A \in \mathbb{R}^{n\times n}$
\begin{enumerate}
    \item The eigenvalues $\lambda$ of A are the roots of the characteristic olynomial $P_A(\lambda) = \det(A-\lambda I)$
    \item The $\lambda$-eigenvectors $\vec{x}$ are the non zero solutions to the homogeneous system $(A - \lambda I)\vec{x}$
\end{enumerate}

\paragraph{Example}
Find the eigenvector related to each eigenvalue:
\begin{align*}
    A = \begin{bmatrix}
        3 & 5 \\
        1 & -1
    \end{bmatrix}, \lambda_1 = 4, \lambda_2 = -2\\
       (A - \lambda_1 I) = \begin{bmatrix}
        3 - 4 & 5 \\
        1 & -1 - 4
        \end{bmatrix} = \begin{bmatrix}
            5 \\
            1
        \end{bmatrix}\\
        (A - \lambda_2 I) = \begin{bmatrix}
            5 & 5 \\
            1 & 1
        \end{bmatrix} \vec{x} = \vec{0} \\
        = \begin{bmatrix}
            1 & 1 \\
            0 & 0
        \end{bmatrix}\vec{x} \implies \vec{x} = \begin{bmatrix}
            -t \\
            t
        \end{bmatrix} = t \begin{bmatrix}
            -1 \\
            1
        \end{bmatrix} \\
        \therefore \begin{bmatrix}
            -1 \\
            1
        \end{bmatrix} \text{ is a basic eigenvector of A corresponding to } \lambda_2 = -2
\end{align*}


\subsubsection*{Transposition}
\paragraph*{Theorem}
$A \in \mathbb{R}^{n\times n}$ and $A^T$ have the same characteristic polynomial.
Therefore they have the same eigenvalues. $A\vec{x} = \lambda \vec{x} \land A^T\vec{x} = \lambda \vec{x}$

\paragraph*{Example}
Find the a. characteristic polynomial b. eigenvalues c. eigenvectors for $A = \begin{bmatrix}
    2 & 0 & 0 \\
    1 & 2 & -1 \\
    1 & 3 & -2
\end{bmatrix}$
\begin{align*}
    P_A(\lambda) = \det(A - \lambda I) = \det\begin{vmatrix}
        2 - \lambda & 0 & 0 \\
        1 & 2 - \lambda & -1 \\
        1 & 3 & -2 - \lambda
    \end{vmatrix} = (2 - \lambda) \det\begin{vmatrix}
        2 - \lambda & -1 \\
        3 & -2 - \lambda
    \end{vmatrix} = (2 - \lambda) ((2 - \lambda)(-2 - \lambda) + 3) = (2 - \lambda)(\lambda^2 + 1) = (2-\lambda)(\lambda - 1)(\lambda + 1)\\
    \therefore \lambda_1 = 2, \lambda_2 = 1, \lambda_3 = -1\\
    (A - 2I) = \begin{bmatrix}
    0 & 0 & 0 \\
    1 & 0 & -1 \\
    1 & 3 & -4
    \end{bmatrix} \xmapsto{RREF} = \begin{bmatrix}
    1 & 0 & -1 \\
    0 & 1 & -1 \\
    0 & 0 & 0
\end{bmatrix} \implies \vec{x} = t \begin{bmatrix}
    1 \\
    1 \\
    1
\end{bmatrix} \text{ is a basic eigenvector corresponding to } \lambda_1 = 2 \\
    (A - I) = \begin{bmatrix}
    1 & 0 & 0 \\
    1 & 1 & -1 \\
    1 & 3 & -3
    \end{bmatrix} \xmapsto{RREF} = \begin{bmatrix}
    1 & 0 & 0 \\
    0 & 1 & -1 \\
    0 & 0 & 0
\end{bmatrix} \implies \vec{x} = t \begin{bmatrix}
    0 \\
    1 \\
    1
\end{bmatrix} \text{ is a basic eigenvector corresponding to } \lambda_2 = 1 \\
    (A + I) = \begin{bmatrix}
    3 & 0 & 0 \\
    1 & 3 & -1 \\
    1 & 3 & -1
    \end{bmatrix} \xmapsto{RREF} = \begin{bmatrix}
    1 & 0 & 0 \\
    0 & 1 & -\frac{1}{3} \\
    0 & 0 & 0
\end{bmatrix} \implies \vec{x} = t \begin{bmatrix}
    0 \\
    \frac{1}{3} \\
    1
\end{bmatrix} \text{ is a basic eigenvector corresponding to } \lambda_3 = -1
\end{align*}

S